
\chapter{ISO 8583: формат сообщений}\label{ch:iso8583}

В~процессинге авторизационный запрос на~покупку 1\,250~рублей 50~копеек занимает 153~байта. Чтобы прочитать дамп как~рассказ о~транзакции, нужно знать грамматику ISO~8583: где заканчивается тип сообщения, где начинаются значения, как~искать конкретное поле. Глава учит это делать на~одном MTI~0100 и~его ответе MTI~0110~--- том~же жизненном цикле, что в~гл.~\ref{ch:tx-lifecycle}, только в~байтах.

\needspace{4\baselineskip}
\section{Структура сообщения}\label{sec:iso8583-structure}

Сообщение целиком~--- на~рис.~\ref{fig:mti-0100-anatomy}. Сверху приведены байты в~hex-виде с~подсветкой по~полям; снизу~--- что несёт каждое поле, где оно начинается и~какой у~него декодированный смысл. Поля переменной длины (LLVAR) выделены другим цветом; первые два байта каждого LLVAR~--- префикс длины~\texttt{LL}~--- показаны насыщеннее остального значения.

\begin{figure}[tbp]
  \centering
  {\renewcommand{\figuremaxheight}{.92\textheight}%
  \includefiguresvg[width=\linewidth]{assets/figures/ch06-iso8583/mti-0100-anatomy}}
  \caption{Разбор MTI~0100 по~байтам: hex-дамп сверху и~таблица полей снизу.}\label{fig:mti-0100-anatomy}
\end{figure}

Грамматика ISO~8583 одинакова для~всех сообщений и~состоит из~трёх частей, идущих строго в~одном порядке:

\begin{enumerate}
  \item \textbf{MTI} (Message Type Indicator)~--- четыре первых байта (\texttt{30 31 30 30}, в~ASCII даёт~\texttt{0100}). Отвечают на~вопрос, \emph{что это за~сообщение}: авторизация, финансовое сообщение, отмена, сетевое управление.
  \item \textbf{Bitmap}~--- следующие восемь байт (\texttt{72 3C 04 81 20 C0 80 00}). Битовая карта, отвечающая на~вопрос, \emph{какие поля присутствуют дальше}.
  \item \textbf{Элементы данных} (data elements, DE)~--- сами значения: в~нашем сообщении пятнадцать полей: PAN (Primary Account Number; разбор~--- в~гл.~\ref{ch:tx-lifecycle}), сумма, STAN (Systems Trace Audit Number, номер трассировки), код обработки, идентификаторы терминала и~ТСП, временные метки и~ещё несколько признаков транзакции.
\end{enumerate}

\highlight{ISO~8583 устроен иначе, чем самоописывающиеся форматы вроде JSON или~XML. Сначала идёт тип операции, затем bitmap с~картой присутствия, затем значения подряд в~порядке возрастания номеров}. Тег и~длина внутри каждого поля не~передаются.

В~примере все числовые поля закодированы в~ASCII~--- каждая цифра занимает один байт. Реальные IFS могут паковать числа в~BCD (Binary Coded Decimal), вдвое сжимая числовые поля; подробности кодировки, ответ эмитента MTI~0110 и~отмены~--- в~разделе~\ref{sec:iso8583-hex-dump}.

\needspace{4\baselineskip}
\section{MTI и~стадия транзакции}\label{sec:iso8583-mti}

MTI~--- первые четыре цифры сообщения; каждая цифра кодирует своё измерение: версию стандарта, класс операции, функцию в~потоке и~инициатора пары. Раскладка для~MTI~0100~--- на~рис.~\ref{fig:mti-breakdown}.

\begin{figure}[tbp]
  \centering
  \includefiguresvg[width=.95\linewidth]{assets/figures/ch06-iso8583/mti-breakdown}
  \caption{MTI~0100 по~четырём позициям; жёлтым подсвечены значения примера.}\label{fig:mti-breakdown}
\end{figure}

В~паре запрос/ответ меняется только третья цифра~--- функция; класс, версия и~инициатор постоянны. Массовые карточные интеграции работают на~версии~\texttt{0} (ISO~8583:1987); редакции 1993 и~2003~годов в~промышленной эксплуатации не~прижились. 4-я~цифра отмечает сторону, начавшую обмен, а~не отправителя конкретного сообщения: в~ответе значение инициатора сохраняется.

\highlight{Request идёт end-to-end, advice и~notification~--- point-to-point}. Эквайер отправляет запрос и~ждёт решения эмитента; advice и~notification на~уровне приложения не~отвергаются, гарантия доставки достигается повторами с~теми~же ссылочными полями. Повтор кодируется заменой 4-й~цифры на~нечётную: MTI~\texttt{0121}~--- та~же advice, что MTI~\texttt{0120}, но~переотправленная при~недоставке.

Из теоретического пространства десяти тысяч комбинаций большая часть нерабочая. MTI~\texttt{1102} формально означал бы authorization request от~эмитента, чего в~карточной модели не~бывает: запрос на~авторизацию всегда инициирует эквайер. Учебные обзоры раскладывают MTI по~разрядам, но~реальный список MTI задаёт IFS партнёра (Interface Specification, подробно~--- в~разделе~\ref{sec:iso8583-implementations}): стандарт фиксирует общую схему, а~конкретный процессинг или~сеть оставляют от~неё узкий рабочий набор~\cite{iso-8583-2023,fis-iso8583-guide-2023}.

Канонические пары запрос/ответ для~ISO~8583:1987:

\begin{itemize}
  \item MTI~0100 / MTI~0110~--- authorization request / response: авторизация без~списания;
  \item MTI~0120 / MTI~0130~--- authorization advice / advice response: уведомление об~операции, прошедшей через \textit{stand-in} или~офлайн;
  \item MTI~0200 / MTI~0210~--- financial request / response: авторизация и~финансовое представление одним сообщением;
  \item MTI~0220 / MTI~0230~--- financial advice / advice response: офлайновая или~force post;
  \item MTI~0320 / MTI~0330~--- file action advice / response: обновление справочников и~BIN-таблиц;
  \item MTI~0400 / MTI~0410~--- reversal request / response: отмена в~реальном времени;
  \item MTI~0420 / MTI~0430~--- reversal advice / response: отложенная отмена;
  \item MTI~0500 / MTI~0510~--- reconciliation request / response: сверка итогов;
  \item MTI~0600 / MTI~0610~--- административные сообщения;
  \item MTI~0800 / MTI~0810~--- network management: echo test, обмен ключами, sign-on/off.
\end{itemize}

В~паре MTI~0100/MTI~0110 эмитент принял решение об~авторизации, а~деньги в~расчётном смысле ещё не~двинулись; перевод в~клиринг произойдёт отдельным сообщением, доставляемым клиринговым файлом (BASE~II~TC у~Visa, IPM~MTI~1240 у~Mastercard, собственный профиль ISO~8583 у~НСПК для~карт \enquote{Мир}~--- разбор в~гл.~\ref{ch:visa},~\ref{ch:mastercard}~и~\ref{ch:mir}), и~инициирует его эквайер, не~эмитент.

MTI~0120~--- advice: парный~0130 определён спецификацией, но~многие IFS его не~реализуют, считая доставку самой advice достаточной. Отложенная авторизация (\textit{deferred authorization})~--- когда операция должна пройти быстрее, чем позволяет онлайн-запрос к~эмитенту: транспортные гейты решают за~доли секунды по~локальным правилам, advice MTI~0120 уходит уже после факта~\cite{uspf-transit-2018}. \textit{stand-in} processing (STIP)~--- когда эмитент или~его процессинг недоступны: карточная сеть отвечает на~запрос авторизации от~его имени по~заранее заданным правилам, advice досылается эмитенту после восстановления связи~\cite{visa-stip-2020}.

Пара MTI~0200/MTI~0210~--- финансовое сообщение; в~карточных диалектах именно она реализует одностадийный сценарий с~авторизацией и~финансовым представлением в~одной транзакции.

Отмена идёт через MTI~0400/MTI~0410 в~одних интерфейсах и~через MTI~0420/MTI~0430 в~других. Разница повторяет водораздел request/advice: 0400~--- запрос на~отмену в~реальном времени, на~который эмитент отвечает; 0420~--- одностороннее уведомление об~отмене. Какую пару применять, диктует спецификация интерфейса~\cite{fis-iso8583-guide-2023}.

Служебные пары MTI~0800/MTI~0810 обслуживают \emph{сетевой диалог}: \enquote{sign-on}~--- логический вход в~сеть с~началом сессии и~обновлением ключей; \enquote{echo} (\enquote{echo-test})~--- периодическая проверка живости канала; \enquote{cutover}~--- переключение на~новый набор ключей или~новую версию профиля без~разрыва связи.

\needspace{4\baselineskip}
\section{Bitmap}\label{sec:iso8583-bitmap}

\highlight{После MTI идёт bitmap~--- восемь или~шестнадцать байт, в~которых каждый бит отвечает за~присутствие одного элемента данных}.
Сами значения bitmap не~несёт; он показывает только, какие поля искать дальше в~теле сообщения~\cite{fis-iso8583-guide-2023}.

Поздние форматы вроде TLV (Tag-Length-Value), как в~DE~55 для~EMV (Europay\slash Mastercard\slash Visa, чиповый платёжный стандарт; подробно~--- в~гл.~\ref{ch:emv}), делают сообщение самоописывающимся, и~каждое поле несёт собственный тег и~длину. Bitmap проектировался в~1970--80-х, когда каждый байт в~телекоммуникационном канале имел цену, и~проверить бит маской было дешевле, чем пройти по~тегам.

В~нашем примере основная bitmap~--- \texttt{72 3C 04 81 20 C0 80 00}. Развёрнутая в~биты, она показывает, какие DE присутствуют в~сообщении: бит~N в~bitmap~--- флаг присутствия DE~N. Биты внутри байта нумеруются слева направо, от~старшего (MSB, Most Significant Bit~--- бит с~наибольшим весом, $2^7 = 128$ в~восьмибитном байте) к~младшему (LSB, Least Significant Bit~--- бит с~наименьшим весом, $2^0 = 1$). В~нашей таблице бит~1 каждого байта~--- его MSB; разворот по~ячейкам~--- на~рис.~\ref{fig:bitmap-presence}.

\begin{figure}[tbp]
  \centering
  \includefiguresvg[width=.85\linewidth]{assets/figures/ch06-iso8583/bitmap-presence}
  \caption{Bitmap MTI~0100, развёрнутая в~64~ячейки: заполненная~--- DE присутствует.}\label{fig:bitmap-presence}
\end{figure}

Первый бит первого байта (\texttt{72} = \texttt{0111\,0010}) равен~0~--- значит, дополнительной (secondary) битовой карты нет и~сообщение укладывается в~первые 64~DE. Остальные единицы дают: DE~2, 3, 4, 7, 11, 12, 13, 14, 22, 25, 32, 35, 41, 42, 49~--- правдоподобный набор для~обычной авторизации, в~котором есть идентификация карты, сумма, временной контекст, ТСП и~способ ввода.

Если~бы старший бит \texttt{72} стоял в~единице, после первой bitmap шли~бы ещё восемь байт~--- дополнительная (secondary) bitmap для~DE~65\dots128. Старший бит уже её первого байта аналогично включает третью (tertiary) bitmap на~DE~129\dots192: так читатель встречает три яруса присутствия~--- основная (primary), дополнительная (secondary), третьего уровня (tertiary). Secondary в~реальных профилях обычна (Visa BaseI и~Mastercard CIS возят её рутинно); tertiary публично описана только в~Visa BaseI и~встречается редко, обычно под~специализированные служебные поля.

Декодер принимает байты сообщения и~возвращает MTI и~список номеров присутствующих DE; чтение самих значений требует длины и~формата каждого поля из~IFS партнёра.

\codefile{Python}{samples/ch06-iso8583-bitmap/bitmap.py}

\practicenote{%
  При~разборе инцидента сначала смотрите на~bitmap. Он быстро
  показывает, есть ли в~сообщении DE~55, Track~2, DE~38 и~DE~39. Часто
  этого хватает, чтобы снять половину вопросов ещё до~полного чтения дампа.

}

Декодер выше намеренно ограничивается primary и~secondary bitmap и~возвращает голый список номеров полей: MTI~--- шаг первый, bitmap~--- второй. Реальный разбор сообщения добавляет к~этому третий шаг: для~каждого номера DE из~списка взять формат и~длину из~IFS партнёра, отрезать нужное количество байт от~тела сообщения и~декодировать значение.

У~полей фиксированной длины она зашита в~спецификации; у~LLVAR и~LLLVAR префикс длины идёт первым внутри самого поля. \texttt{LL}~--- двузначный префикс длины, \texttt{LLL}~--- трёхзначный, и~то, сколько байт они занимают, зависит от~кодировки: в~ASCII LL занимает 2~байта, LLL~--- 3~байта; в~BCD те~же цифры пакуются вдвое плотнее, LL умещается в~1~байте, LLL~--- в~2 (со~старшим нулём для~выравнивания по~границе байта). Кодировка длины может не~совпадать с~кодировкой значения и~оговаривается отдельной строкой IFS, прочесть её надо до~разбора поля.

\needspace{4\baselineskip}
\section{DE~--- элементы данных}\label{sec:iso8583-data-elements}

Для~MTI~0100/MTI~0110 работающие поля распадаются по~пяти смысловым группам; полный каталог DE с~форматами и~длинами берётся из~сетевой спецификации и~IFS партнёра~\cite{fis-iso8583-guide-2023}.

\subsection*{Идентификация и~маршрутизация}

DE~2 (PAN) говорит, чья карта. DE~35 (Track~2 Data) повторяет тот же PAN в~более богатой форме и~укладывает его вместе со~сроком действия, сервисным кодом и~дискреционными данными в~одну строку.
DE~32~--- идентификатор эквайера или~банка-инициатора. DE~37 (RRN, Retrieval Reference Number~--- сквозной ссылочный номер операции) проходит через запрос, ответ и~все последующие сообщения, когда операцию приходится искать в~отчётах, диспутах и~операторских инструментах~\cite{fis-iso8583-guide-2023}.

\subsection*{Деньги и~время}

DE~3 (Processing Code) кодирует тип операции: покупка, возврат, снятие наличных, запрос баланса.
DE~4 хранит сумму с~фиксированным масштабом, который определяется кодом валюты из~DE~49 по~таблице экспонентов ISO~4217.
DE~7 фиксирует сетевое время передачи сообщения. DE~11~--- STAN (Systems Trace Audit Number), шестизначный номер трассировки, по~которому процессинг связывает запрос с~ответом. DE~12 и~DE~13~--- локальное время и~дата терминала. DE~49~--- код валюты.

STAN уникален только в~пределах сетевого времени из~DE~7; локальное время терминала в~споре с~эмитентом ориентиром не~служит.

\subsection*{Контекст точки приёма}

DE~22 (POS Entry Mode) показывает, как получены данные карты: вручную, с~магнитной полосы, с~чипа, бесконтактно, в~интернет-операции. DE~25 (POS Condition Code) добавляет условия операции. DE~41 и~DE~42~--- идентификаторы терминала и~ТСП~--- привязывают транзакцию к~точке приёма.

Одни и~те же PAN и~сумма получают у~эмитента разную оценку~--- чиповая покупка в~супермаркете и~ручной ввод реквизитов на~сайте идут по~разным шкалам риска.

\subsection*{Результат авторизации}

DE~39 (Response Code)~--- главное поле ответа: одобрена операция или~нет. DE~38 (идентификатор авторизации) появляется при~одобрении и~связывает авторизацию с~последующим клирингом.

В~запросе этой группы нет: решение ещё не~принято. Чтение MTI~0110 начинается с~DE~39.

\subsection*{Чиповые данные}

DE~55~--- контейнер TLV для~данных EMV, в~котором содержатся криптограммы, TVR (Terminal Verification Results, результаты проверок терминала), AID (Application Identifier, идентификатор платёжного приложения карты), ATC (Application Transaction Counter, счётчик операций приложения) и~другие теги, которыми терминал доказывает эмитенту, что операция пришла с~чипа или~ядра. Подробный разбор тегов~--- в~гл.~\ref{ch:emv}~\cite{emvco-specs}.

DE~55 в~примере умышленно опущен: чиповый хвост в~реальной авторизации значительно увеличивает размер сообщения, и~без него проще разделить грамматику ISO~8583 и~детали EMV.

\needspace{4\baselineskip}
\section{Ответ, отмена и~нюансы кодирования}\label{sec:iso8583-hex-dump}

В~запросе на~рис.~\ref{fig:mti-0100-anatomy} эквайер просит одобрить бесконтактную покупку на~сумму 1\,250~рублей 50~копеек по~карте \texttt{2200121234560892} в~терминале \texttt{TERM0001} у~ТСП \texttt{MERCHANT0000001}.

\subsection*{MTI~0110~--- ответ эмитента}

MTI~0110 короче запроса: PAN и~Track~2 не~нужны, зато появляются DE~39 (response code) и, при~одобрении, DE~38 (authorization ID response). STAN, сумма, дата и~идентификатор терминала повторяются ради сопоставления ответа с~запросом~\cite{fis-iso8583-guide-2023}.

В~нашем сообщении ответный bitmap, например, \texttt{32 38 00 00 0E C0 80 00}~--- DE~3, DE~4, DE~7, DE~11, DE~12, DE~13, DE~37, DE~38, DE~39, DE~41, DE~42, DE~49. DE~37 несёт RRN~--- сквозной ссылочный номер для~сопоставления в~клиринге; DE~38~--- код авторизации (в~нашем примере \texttt{A1B2C3}); DE~39 фиксирует решение эмитента: \texttt{00}~--- одобрение, \texttt{51}~--- недостаток средств~\cite{fis-iso8583-guide-2023}.

\subsection*{Отмена (reversal)~--- MTI~0400 и~MTI~0420}

Одни интерфейсы оформляют отмену как MTI~0400/MTI~0410, другие как MTI~0420/MTI~0430. FIS/Worldpay разделяет варианты request и~advice~\cite{fis-iso8583-guide-2023}. Смысл общий~--- \enquote{считать предыдущее сообщение несостоявшимся} или~\enquote{скорректировать его последствия}. Возврат денег клиенту~--- отдельная операция; различия reversal, void и~refund подробно разобраны в~разделе~\ref{sec:tx-cancellation}.

Сообщение об~отмене должно однозначно сослаться на~исходную операцию. Для~этого используются тот же STAN (DE~11), тот же временной контекст (DE~7), а~в~ряде профилей~--- специальные \enquote{original data elements} (поля, в~которые упаковываются ссылки на~исходную операцию: её MTI, STAN, дата и~время) вроде~DE~90. Официальные профили прямо описывают DE~90 как способ передать эмитенту значения исходной операции~\cite{fis-iso8583-guide-2023}.

\subsection*{Чего не~видно в~hex}

Один и~тот же байт значит разное в~разных профилях, а~внутри одного поля прячутся собственные подполя со~своей структурой.

\textbf{POS Entry Mode (DE~22).} В~нашем сообщении DE~22=\texttt{071}. В~этом профиле код описывает бесконтактный чиповый сценарий с~терминалом, способным принять PIN (третий разряд кодирует PIN entry capability терминала; факт верификации PIN в~операции через DE~22 не~передаётся). В~другом IFS тот же сценарий может кодироваться, например, \texttt{05F}~\cite{fis-iso8583-guide-2023}. Код, работающий с~одним процессингом, у~другого партнёра молча падает~--- такие диалектные расхождения надо ловить на~тестовых транзакциях, в~промышленной эксплуатации (\textit{production}) они обходятся дороже.

\textbf{Сумма с~неявным масштабом (DE~4 и~DE~49).} DE~4 хранит сумму без~десятичной точки. Масштаб задаётся кодом валюты из~DE~49 по~таблице экспонентов ISO~4217: для~рубля и~евро~--- две минорные единицы, для~иены~--- ноль. В~нашем сообщении DE~4=\texttt{000000125050} и~DE~49=\texttt{643} (рубль), что и~даёт 1\,250~рублей 50~копеек.

\textbf{Внутренняя структура Track~2 (DE~35).} Track~2 несёт несколько подполей в~одной строке: PAN, разделитель~\texttt{=}, срок действия~\texttt{YYMM}, трёхзначный сервисный код, далее дискреционные данные. В~примере: \texttt{2200121234560892=26122011234567890}~--- PAN \texttt{2200121234560892}, срок \texttt{2612} (декабрь~2026), сервисный код~\texttt{201}, дальше идут дискреционные данные.

\textbf{Кодировка чисел: ASCII или~BCD.} В~разобранном примере все числовые поля закодированы в~ASCII~--- каждая цифра занимает один байт. В~реальных IFS цифры могут идти в~упакованном числовом формате BCD (Binary Coded Decimal): две цифры в~одном байте, числовые поля вдвое короче. ASCII удобнее при~отладке в~логах, BCD типичен для~промышленной эксплуатации, где каждый лишний байт стоил денег на~телеком-канале и~времени обработки на~железе 1980-х. Прежде чем читать дамп вручную, надо узнать кодировку из~IFS~--- иначе ошибка на~этом этапе делает бессмысленным весь дальнейший разбор~\cite{fis-iso8583-guide-2023}.

\needspace{4\baselineskip}
\section{IFS поверх ISO~8583}\label{sec:iso8583-implementations}

ISO~8583 не~описывает сетевой протокол целиком. Стандарт определяет \emph{грамматику сообщения}; в~официальной аннотации ISO оговорено, что метод расчёта в~предмет стандарта не~входит~\cite{iso-8583-2023,fis-iso8583-guide-2023}.

\textbf{Транспорт и~фрейминг.} Способ доставки~--- по~TCP, через VPN, с~каким заголовком и~префиксом длины~--- стандартом не~предписан. Двух- или~четырёхбайтовый префикс длины, отдельный application header, выбор между ASCII и~EBCDIC (Extended Binary Coded Decimal Interchange Code~--- 8-битная кодировка IBM, исторически используемая на~мейнфреймах процессингов)~--- всё это соглашение конкретного процессинга.
Транспортный нейтралитет был проектным выбором: стандарт 1987~года должен был одинаково работать на~X.25, SNA и~любом будущем протоколе. Цена этого решения в~том, что каждый IFS фиксирует транспорт сам и~при смене партнёра ничего из~этого не~переносится.

\textbf{Криптография и~ключи.} PIN-блок (зашифрованный блок с~PIN-кодом фиксированной длины; форматы~--- в~разделе~\ref{sec:crypto-pin-block}) едет в~DE~52, MAC (Message Authentication Code, код аутентификации сообщения; алгоритмы CBC-MAC и~CMAC, ключи и~состав строки~--- в~разделе~\ref{sec:crypto-symmetric})~--- в~DE~64 или~DE~128. Алгоритмы, формат PIN-блока, схема ключей и~порядок их обмена задаются сетью или~спецификацией конкретного коммутатора; сам ISO~8583 их не~описывает.
\highlight{Само наличие DE~64 или~DE~128 не~объясняет MAC}: для~проверки
нужны ключи, алгоритм, padding, состав строки и~транспортная
спецификация~--- обвязка вокруг ISO~8583, лежащая вне его
базовой грамматики.

\textbf{Состояние сессии.} Sign-on, echo test, cutover, управление ключами и~повторная отправка опираются на~MTI~0800/MTI~0810 и~на~то, как конкретная система понимает таймауты, идемпотентность (гарантию, что повторная отправка того же сообщения не~приведёт к~двойному списанию или~двойной отмене) и~повтор.

\textbf{Правила обработки.} Что считать дубликатом, когда отправлять сообщение отмены, как долго ждать ответ, какие поля обязательны для~чиповой или~интернет-операции~--- всё это прикладные правила интерфейса, лежащие вне синтаксиса сообщения.

В~одной системе всё кодируется в~ASCII, в~другой рядом живут ASCII, BCD и~бинарные поля. Один партнёр считает DE~32 обязательным, другой не~принимает его вовсе; одному достаточно DE~35, другому нужен ещё DE~2, третьему~--- набор частных полей под~конкретный сценарий. Где-то дополнительные данные раскрываются через DE~48, DE~60, DE~61, DE~63 или~DE~124, где-то собираются в~собственные подполя. Аналогично расходятся коды ответа, жизненный цикл сообщений отмены, cutover, транспортный заголовок, sign-on и~ключевой обмен. Какие поля обязательны для~конкретного канала, как кодируется сумма, где искать частные подполя и~каким должен быть ответ на~MTI~0800~--- отвечает только IFS партнёра~\cite{fis-iso8583-guide-2023}.

Детали диалектов Visa, Mastercard и~НСПК~--- в~их сетевых спецификациях, как и~полный справочник элементов данных и~дампы MTI~0200, MTI~0420 и~сетевых служебных сообщений.
